\section{L'idea: unire logica e intelligenza artificiale}

Come abbiamo visto nel capitolo precedente, il vero limite dei moderni sistemi di
riparazione automatica del codice non è la mancanza di ``creatività'' nel generare
soluzioni, ma la mancanza di una \textbf{guida rigorosa}. Quando lasciamo che un Large Language
Model lavori da solo, senza vincoli precisi, finisce per inventare parti di codice inesistenti
(allucinazioni) o per produrre soluzioni che funzionano solo in apparenza (overfitting).
All'estremo opposto, gli strumenti puramente logici e formali sono estremamente precisi
nell'individuare \textit{dove} si trova l'errore, ma non hanno la flessibilità necessaria
per scrivere il codice che lo risolve in un contesto reale.

L'idea alla base di \textbf{HybridRepair} nasce proprio da qui: il mondo logico/simbolico
e quello neurale/LLM non sono in competizione, ma \textbf{complementari}. 
per la diagnosi si impiega \textbf{LogicFL}, uno strumento 
preesistente, ma il contributo di questo lavoro è la 
progettazione dell'\textbf{architettura di integrazione} che
trasforma il suo referto formale in una guida strutturata per
la sintesi LLM, con una validazione che va oltre i semplici
test automatici.

\subsection{La filosofia: separare diagnosi e sintesi}

Il principio fondamentale di HybridRepair è una netta \textbf{separazione dei compiti}:

\begin{quote}
\textbf{L'analisi logica si occupa di trovare la causa del problema (diagnosi) mentre l'LLM si occupa di scrivere il codice che lo risolve (sintesi).}
\end{quote}

\textbf{La diagnosi.} Capire da dove arriva un valore \texttt{null}, seguire il
percorso di un errore attraverso più funzioni e individuare il punto esatto in cui il
programma va in crash richiede un rigore logico e deterministico. Sono garanzie che solo
la programmazione logica (come Prolog) può offrire.

\textbf{La sintesi.} Una volta che sappiamo con certezza \textit{cosa}
correggere, il problema diventa \textit{come} scriverlo in codice Java pulito, che rispetti lo
stile del progetto e superi i test automatici. Qui gli LLM sono insuperabili: comprendono
il contesto e generano codice flessibile con grande efficacia.

\subsection{L'approccio ibrido in pratica: l'esempio Csv-4}

Per chiarire il funzionamento, analizziamo un caso reale: il bug \textbf{Csv-4}, un errore
di tipo \texttt{NullPointerException} (NPE) tratto dal benchmark Defects4J e situato
all'interno di Apache Commons CSV, nella classe \texttt{CSVParser}.

Questo esempio è interessante per un motivo controintuitivo: è un caso in cui \emph{sia}
un sistema basato solo su LLM \emph{sia} l'approccio ibrido arrivano alla patch corretta.
Proprio per questo è utile come esempio motivazionale: ci permette di spostare l'attenzione
dal banale ``chi supera i test'' al punto che davvero conta, cioè \textbf{con quali garanzie}
e attraverso quale processo si arriva a una soluzione che non solo ``sembra giusta'', ma
\emph{è} giusta. Come vedremo, è esattamente qui che emerge il contributo distintivo
dell'approccio ibrido: non tanto nel \emph{generare} la correzione, cosa che oggi un buon
LLM sa fare, quanto nel \textbf{produrre da se la diagnosi} che rende quella correzione
più affidabile. È la concretizzazione del principio ``separare diagnosi e sintesi'' enunciato sopra.

\subsubsection{Il problema}

Il metodo \texttt{getHeaderMap()} restituisce una \emph{copia} della mappa che associa i
nomi delle colonne alla loro posizione. Il suo Javadoc, però, si limita a dire ``returns a copy
of the header map'': non specifica cosa accada quando un header non è definito. Quella parte del
contratto non è documentata: è fissata dall'\emph{oracolo di test} \texttt{testNoHeaderMap()},
che con un \texttt{assertNull} richiede che, in assenza di header, il metodo restituisca
\texttt{null} e non una mappa vuota. È una specifica \emph{comportamentale}, confermata anche
dalla correzione applicata dagli sviluppatori di Apache Commons CSV, secondo cui \texttt{null}
(``header assente'') e \texttt{\{\}} (``header presente, zero colonne'') non sono intercambiabili.

Nel codice originale, quando manca l'header, il campo \texttt{this.headerMap} resta
\texttt{null}; alla riga 288 il metodo tenta comunque di costruirne una copia
(\texttt{new LinkedHashMap<>(null)}) e il risultato è il crash:

\begin{lstlisting}[language=Java]
public Map<String, Integer> getHeaderMap() {
    return new LinkedHashMap<String, Integer>(this.headerMap); // NPE se headerMap e' null
}
\end{lstlisting}

Il \texttt{null} non nasce qui: nasce nel metodo \texttt{initializeHeader()}, dove la variabile
\texttt{hdrMap} viene inizializzata a \texttt{null} e restituita immutata quando non
c'è un header; da lì viene assegnata al campo nel costruttore e infine
provoca il crash. È una catena causale che attraversa più metodi.

\subsubsection{La trappola: una patch che ``sembra giusta''}

C'è una correzione che viene immediatamente in mente e che sembra perfetta: se manca l'header,
restituiamo una \textbf{mappa vuota}.

\begin{lstlisting}[language=Java]
public Map<String, Integer> getHeaderMap() {
    // "Niente header? Restituiamo una mappa vuota": il crash sparisce...
    return this.headerMap == null ? new LinkedHashMap<>() : new LinkedHashMap<>(this.headerMap);
}
\end{lstlisting}

Compila, elimina l'NPE e a prima vista è del tutto ragionevole. In realtà
\textbf{viola il contratto}: l'API deve restituire \texttt{null}, non \texttt{\{\}}. Il test lo
smaschera subito:

\begin{lstlisting}
java.lang.AssertionError: expected null, but was:<{}>
    at CSVParserTest.testNoHeaderMap(CSVParserTest.java:670)
\end{lstlisting}

Questa è la classica insidia di un approccio puramente ``code-centric'': una patch che toglie il
sintomo (il crash) ma cambia in silenzio il significato del metodo.\footnote{Non è un esempio
costruito a tavolino: la patch alla mappa vuota è esattamente la \emph{prima} soluzione proposta
dall'LLM all'interno di HybridRepair, prima che il loop di validazione la respingesse (cfr.
Step 3).} Il punto chiave è che, per evitarla, non serve un'IA ``più intelligente'': serve che il
sistema sappia con precisione \textit{qual è il comportamento atteso}.

\subsubsection{Come la affronta VibeRepair (l'IA che ragiona sulla specifica)}

VibeRepair \cite{zhu2026specification} è uno dei migliori sistemi APR su Defects4J basati su LLM
del 2026. Non è un LLM ``ingenuo'' che riscrive il codice in un colpo solo: è un approccio
\textbf{specification-centric}. Invece di modificare direttamente il codice, traduce il metodo
difettoso in una \emph{specifica del comportamento}, ne inferisce l'intento corretto, ripara la
specifica e solo allora rigenera il codice; inoltre usa il feedback dei test falliti per
raffinare la propria ipotesi. Su Csv-4 questo processo lo porta alla patch corretta:

\begin{lstlisting}[language=Java]
public Map<String, Integer> getHeaderMap() {
    if (this.headerMap == null) {
        return null;
    }
    return new LinkedHashMap<String, Integer>(this.headerMap);
}
\end{lstlisting}

VibeRepair, quindi, \textbf{non cade nella trappola}: evita sia il crash sia la mappa vuota. Vale
però la pena notare \emph{su cosa} poggia questo successo. Primo: come la maggior parte dei
sistemi LLM su Defects4J, VibeRepair riceve già la zona difettosa su cui lavorare (la classica
assunzione di ``fault localization perfetta''): non individua da sé il guasto. 
Secondo: la sua diagnosi è una specifica \emph{inferita} dal solo metodo difettoso e raffinata sul
feedback dei test falliti, non dalla conoscenza dell'intero codice, un'astrazione che può catturare bene
l'intento, come in questo caso, oppure mancarlo. La correttezza, in altre parole, dipende
dalla \textit{bontà di un'inferenza}, non da una garanzia formale.

\subsubsection{Come la affronta HybridRepair (diagnosi formale + sintesi guidata)}

L'approccio ibrido raggiunge la stessa patch corretta, ma per una strada diversa, in tre passi:

\begin{enumerate}
    \item \textbf{Step 1 (LogicFL, diagnosi).} È il passo che fa la differenza. Senza ricevere
    alcun suggerimento su \emph{dove} si trovi il bug, il motore logico esegue il test fallito e
    ricostruisce \emph{autonomamente} l'intera catena dell'errore (non solo la riga del crash):
\begin{lstlisting}[language=Prolog]
npe(line(288), null_expr(headerMap)).
caused_by(npe(line(288)), null_literal(line(325))).     % Map<...> hdrMap = null;
caused_by(npe(line(288)), return(line(351), hdrMap)).   % initializeHeader() ritorna hdrMap
caused_by(npe(line(288)), assign(line(244), headerMap)).% this.headerMap = initializeHeader()
\end{lstlisting}

    \textbf{Questa catena non è lo stack trace.} Lo stack
    trace dell'NPE contiene solo il punto di crash, la
    riga 288, più due frame interni della JDK
    (\texttt{LinkedHashMap.<init>}) e il test chiamante,
    ma \emph{non} le righe 325, 351 e 244: quelle vengono
    eseguite durante la costruzione dell'oggetto, sono già
    ritornate quando l'eccezione scatta e perciò non
    figurano sulla pila. Lo stack trace dice \emph{dove} il
    programma è morto; la catena di LogicFL dice \emph{da
    dove arriva il \texttt{null}} e come si propaga, è
    un'analisi di flusso dei dati.

    Quindi all'LLM non viene detto genericamente ``c'è un crash da qualche parte''. Riceve un referto
    preciso: \textit{il \texttt{null} nasce alla riga 325, si propaga attraverso le righe 351 e
    244, e provoca il crash alla riga 288}. È esattamente il passo che un sistema puramente LLM
    \textbf{non compie}: dove VibeRepair dà per scontato di sapere già quale metodo toccare, qui
    la localizzazione del guasto e la sua causa non sono un'assunzione di partenza, ma un
    \textbf{risultato dimostrato} dal motore logico.

    \item \textbf{Step 2 (Estrazione AST, contesto).} \\
    Il sistema isola le sole ${\sim}35$ righe dei due metodi rilevanti
    (cioè \texttt{getHeaderMap} e \texttt{initializeHeader}), scartando le altre
    ${\sim}460$ righe del file. In questo modo l'LLM ha una visuale pulita e mirata sul
    punto da correggere.

    \item \textbf{Step 3 (Sintesi LLM, con paletti e validazione).} \\
    L'LLM riceve la causa esatta, un contesto pulito e una regola esplicita nel prompt:
    \emph{non far sparire il crash con un oggetto fittizio; superare i test non basta}.
    Ciononostante, la sua prima proposta restituisce di fatto una mappa vuota; è il loop di
    validazione, rieseguendo la suite, a respingerla con \texttt{expected null, but was:<\{\}>}.
    È il test a fissare il comportamento atteso; l'LLM legge la traccia e corregge in \texttt{null}:
\begin{lstlisting}[language=Java]
public Map<String, Integer> getHeaderMap() {
    return this.headerMap == null ? null : new LinkedHashMap<>(this.headerMap);
}
\end{lstlisting}
    Tutti i 48 test passano e il controllo semantico conferma che il contratto è rispettato.
\end{enumerate}

Questa patch ottiene esattamente lo stesso risultato della correzione applicata nella realtà
dagli sviluppatori di Apache Commons CSV. La differenza, in definitiva, sta negli \textbf{input}: 
VibeRepair esige che il guasto sia
\emph{già localizzato}; HybridRepair parte dagli stessi test e dal codice e la localizzazione,
con la catena causale, la \emph{produce}. Richiede cioè meno informazioni in ingresso e ne
restituisce di più, per dimostrazione, non per inferenza, e in modo riproducibile.

\subsubsection{Il confronto}

\begin{table}[H]
\centering
\caption{Confronto tra VibeRepair e HybridRepair sul bug Csv-4}
\vspace{2mm}
\small
\begin{tabularx}{\textwidth}{l >{\raggedright\arraybackslash}X >{\raggedright\arraybackslash}X}
\toprule
\textbf{Caratteristica} & \textbf{VibeRepair (Solo IA)} & \textbf{HybridRepair (Logica + IA)} \\
\midrule
Input richiesti & Codice eseguibile + test falliti + \emph{posizione del guasto} (data)
  & Codice eseguibile + test fallito \\
Localizzazione del guasto & Input (assunta nota: FL perfetta)
  & Output (calcolata da LogicFL) \\
Causa dell'errore & Inferita come specifica (astrazione)
  & Dimostrata come catena di flusso dati (4 livelli) \\
Cosa vede l'LLM & Specifica di comportamento inferita (+ test falliti)
  & Referto causale + i soli 2 metodi rilevanti \\
Difesa dalla mappa vuota & Raffinamento della specifica sui test
  & Loop di validazione sui test \\
Patch finale & \texttt{return null} (corretta) & \texttt{return null} (corretta) \\
Supera i test? & Sì & Sì \\
È davvero corretto? & Sì & Sì \\
\bottomrule
\end{tabularx}
\end{table}

A differenza dell'esempio ``da manuale'' in cui l'IA fallisce, qui entrambi i sistemi
\emph{ben progettati} arrivano alla soluzione giusta, ed è proprio questo a renderlo
istruttivo. Il vero spartiacque non è ``IA contro ibrido'' in termini di test superati, ma la
\textbf{trappola della mappa vuota}: una patch che ``sembra giusta'' e in cui un approccio
ingenuo o puramente code-centric cadrebbe facilmente. Ciò che protegge da quella trappola non è
quanto è ``intelligente'' l'IA, ma la \textbf{qualità delle informazioni} e il rigore del
processo: VibeRepair ricostruisce l'intento \emph{inferendo} una specifica; HybridRepair lo
\emph{garantisce} con una diagnosi formale e autonoma e con una validazione che riesegue i test.
La differenza, in definitiva, è nel \textit{livello di garanzia}, inferenza contro
dimostrazione, localizzazione fornita contro localizzazione autonoma, in modo riproducibile.
